\lecture{12}{2026.06.01}{}

\section{Fast Fourier Transform}

\begin{definition}[Discrete Fourier Transform]
    Consider $n = m$. Using the coefficients obtained from the Discrete Fourier Transform (DFT) of $\{y_j\}_{j=0}^{2m-1}$, the corresponding trigonometric polynomial is given by \[
        S_m(x) = \frac{a_0 + a_m \cos mx}{2} + \sum_{k=1}^{m-1} (a_k \cos kx + b_k \sin kx)
    \]
    where \begin{equation}
        a_k = \frac{1}{m} \sum_{j=0}^{2m-1} y_j \cos k x_j, k = 0, 1, \dots, n \label{eq:ak}
    \end{equation}
    and \begin{equation}
        b_k = \frac{1}{m} \sum_{j=0}^{2m-1} y_j \sin k x_j, k = 1, 2, \dots, n-1 \label{eq:bk}
    \end{equation}
\end{definition}

This calculation requires \begin{center}
    $(2m)^2$ multiplications and $(2m)^2 - 1$ additions.
\end{center}
which is computationally expensive for large $m$. The Fast Fourier Transform (FFT) is an efficient algorithm to compute the DFT, reducing the computational complexity to \begin{center}
    \red{$O(m \log_2 m)$ multiplications and additions.}
\end{center}

\begin{note}
    FFT is named by SIAM as one of the top 10 algorithms in the 20th century.
\end{note}

First, we extend the definition of DFT to a more general form. For $b_k$'s range we use $0 \le k \le n$ instead of $1 \le k \le n-1$. \[
    b_k = \frac{1}{m} \sum_{j=0}^{2m-1} y_j \sin k x_j, k = 0, 1, \dots, m
\]
as \[
    \sin 0 x_j = 0 \mbox{ and } \sin m x_j = \sin \left(m \left( -\pi + \frac{\pi j}{m} \right)\right) = 0
\]
imply \[
    b_0 = b_m = 0
\]

\begin{definition}[Complex Form of DFT]
    Define \begin{equation}
        c_k = \frac{1}{2m} \sum_{j=0}^{2m-1} y_j e^{-\pi i j k /m},\ k = 0, 1, \dots, 2m-1 \label{eq:ck}
    \end{equation}
\end{definition}

We now have \begin{equation*}
\begin{split}
c_k & = \sum_{j=0}^{2m-1} y_j e^{-\pi ijk/m} \\
& = \sum_{j=0}^{2m-1} y_j e^{ik(-\pi j/m + 2\pi)} = \sum_{j=0}^{2m-1} y_j e^{ik(\pi -\pi j/m + \pi)} \\
& = \sum_{j=0}^{2m-1} y_j e^{-ik x_j} e^{ik\pi} = m(a_k - i b_k) (\cos k\pi + i \sin k\pi)\\
& = m (a_k - i b_k) (-1)^k,
\end{split}
\end{equation*}
where in one place we use
\begin{equation*}
  x_j = -\pi + \frac{\pi j}{m},
\end{equation*}
and $a_k$ and $b_k$ are defined in \eqref{eq:ak} and \eqref{eq:bk} respectively. If $c_k$ is defined as in \eqref{eq:ck}, then we can derive the following relationships: \[
    \frac{c_k \cdot (-1)^k}{m} = a_k - i b_k
\]

The above equation implies that \begin{equation*}
  \begin{split}
    a_k & = \textrm{Re}\left(\frac{c_k (-1)^k}{m}\right)\\
    b_k & = -\textrm{Im}\left(\frac{c_k (-1)^k}{m}\right)
  \end{split}
\end{equation*}
where $\textrm{Re}(z)$ and $\textrm{Im}(z)$ denote the real and imaginary parts of a complex number $z$, respectively.

\subsection{Recursive Form of FFT}

When $m = 2$, \[
    x_j = -\pi + \frac{\pi j}{2}, j = 0, 1, 2, 3
\]
\[
    S_2(x) = \frac{a_0 + a_2 \cos 2x}{2} + a_1 \cos x + b_1 \sin x
\]
Using \eqref{eq:ck}, \begin{equation*}
\begin{array}{rrrrrrr}
c_0 = y_0 e^0 & + &  y_1 e^0         & + & y_2e^0        & + & y_3 e^{0}         , \\   
c_1 = y_0 e^0 & + &  y_1 e^{-\pi i/2} & + & y_2e^{-\pi i}  & + & y_3 e^{-3\pi i/2}  , \\   
c_2 = y_0 e^0 & + &  y_1 e^{-\pi i}   & + & y_2e^{-2\pi i} & + & y_3 e^{-3\pi i}    , \\   
c_3 = y_0 e^0 & + &  y_1 e^{-3\pi i/2}& + & y_2e^{-3 \pi i}& + & y_3 e^{-9 \pi i /2}
\end{array}
\end{equation*}

\begin{note}
    There are some other ways to define $c_k$'s. For example, if \[
        c_k = \sum_{j=0}^{2m-1} y_j e^{\red{\pi i j k /m}},\ k = 0, 1, \dots, 2m-1
    \]
    then \begin{equation*}
    \begin{split}
    c_k & = \sum_{j=0}^{2m-1} y_j e^{ik(-\pi + \pi j/m - \pi)}
    = \sum_{j=0}^{2m-1} y_j e^{ikx_j} e^{-i k \pi}\\
    & = m (a_k + i b_k)(\cos k\pi - i \sin k\pi)\\
    & = m (a_k + i b_k)(-1)^k,
    \end{split}
    \end{equation*}
    But, we do not consider this definition. Instead, we use the definition in \eqref{eq:ck} by following the common convention in the literature.
\end{note}

Now we need a systematic way to calculate \[
    c_k, k = 0, 1, \dots, 2m-1
\]
and then find \[
    a_k, b_k, k = 0, 1, \dots, m
\]
We have \[
    Fy = \begin{pmatrix}
        \vdots \\
        c_k \\
        \vdots
    \end{pmatrix}
\]
where \[
    F = \begin{pmatrix}
        1 & 1 & 1& \cdots & 1 \\
        1 & \delta & \delta^2 & \cdots & \delta^{(2m-1)}\\
        1 & \delta^2 & \delta^4 & \cdots & \delta^{2(2m-1)}\\
        &&\vdots &&\\
        1 & \delta^{(2m-1)} & \delta^{2(2m-1)} & \cdots & \delta^{(2m-1)^2}
    \end{pmatrix}
\]
and \[
    \delta = e^{-\pi i / m}
\]

\begin{note}
    We found that \begin{center}
        $e^{-\pi i / m}$ is the $(2m)$-th root of unity.
    \end{center}
    because \[
        \left(e^{-\pi i / m}\right)^{2m} = e^{-2\pi i} = 1
    \]
\end{note}

Now consider the case of $m = 4$ \[
    F = \begin{pmatrix}
        1 & 1 & 1 & 1 & 1 & 1 & 1 & 1 \\
        1 & \delta^1 & \delta^2 & \delta^3 & \delta^4 & \delta^5 & \delta^6 & \delta^7 \\
        1 & \delta^2 & \delta^4 & \delta^6 & 1 & \delta^2 & \delta^4 & \delta^6 \\
        1 & \delta^3 & \delta^6 & \delta^1 & \delta^4 & \delta^7 & \delta^2 & \delta^5 \\
        1 & \delta^4 & 1 & \delta^4 & 1 & \delta^4 & 1 & \delta^4 \\
        1 & \delta^5 & \delta^2 & \delta^7 & \delta^4 & \delta^1 & \delta^6 & \delta^3 \\
        1 & \delta^6 & \delta^4 & \delta^2 & 1 & \delta^6 & \delta^4 & \delta^2 \\
        1 & \delta^7 & \delta^6 & \delta^5 & \delta^4 & \delta^3 & \delta^2 & \delta^1
    \end{pmatrix}
\]
where \[
    \delta = e^{-\pi i / 4}, \ \delta^4 = e^{-\pi i} = -1,\ \delta^8 = e^{-2\pi i} = 1
\]

Then the $F$ can be rewritten as \begin{equation}
    F = \begin{pmatrix}
        1 & 1 & 1 & 1 & 1 & 1 & 1 & 1 \\
        1 & \delta^1 & \delta^2 & \delta^3 & -1 & -\delta^1 & -\delta^2 & -\delta^3 \\
        1 & \delta^2 & -1 & -\delta^2 & 1 & \delta^2 & -1 & -\delta^2 \\
        1 & \delta^3 & -\delta^2 & \delta & -1 & -\delta^3 & \delta^2 & -\delta^1 \\
        1 & -1 & 1 & -1 & 1 & -1 & 1 & -1 \\
        1 & -\delta^1 & \delta^2 & -\delta^3 & -1 & \delta^1 & -\delta^2 & \delta^3 \\
        1 & -\delta^2 & -1 & \delta^2 & 1 & -\delta^2 & -1 & \delta^2 \\
        1 & -\delta^3 & -\delta^2 & -\delta & -1 & \delta^3 & \delta^2 & \delta^1 \\  
    \end{pmatrix} \label{eq:F4}
\end{equation}

Suppose \[
    c = \begin{pmatrix}
        0 & 2 & 4 & 6 & 1 & 3 & 5 & 7
    \end{pmatrix}
\]
is an \red{index permutation} vector.

Then we rearrange the rows of $F$ and the elements of $y$ according to $c$:
\[
    Fy = F(:, c) y(c)
\]
and \[
    F_8(:, c) = \begin{pmatrix}
        F_4 & \Omega_4 F_4 \\
        F_4 & -\Omega_4 F_4
    \end{pmatrix}
\]

We call the current form of $F$ as $F_8$, then \[
    F_4 = \begin{pmatrix}
        1 & 1 & 1 & 1 \\
        1 & \delta^2 & -1 & -\delta^2 \\
        1 & -1 & 1 & -1 \\
        1 & -\delta^2 & -1 & \delta^2
    \end{pmatrix}
\]
and \[
    \Omega_4 = \begin{pmatrix}
        1& & & \\    
        & \delta & & \\
        & & \delta^2 & \\
        & & & \delta^3 \\
    \end{pmatrix}
\]
Now we have \[
    \Omega_4 F_4 = \begin{pmatrix}
        1 & 1 & 1 & 1 \\
        \delta^1 & \delta^3 & -\delta^1 & -\delta^3 \\
        \delta^2 & -\delta^2 & \delta^2 & -\delta^2 \\
        \delta^3 & \delta^1 & -\delta^3 & -\delta^1 \\
    \end{pmatrix}
\]

Then \begin{align*}
    Fy = F(:, c) y(c) &= \begin{pmatrix}
        F_4 & \Omega_4 F_4 \\
        F_4 & -\Omega_4 F_4
    \end{pmatrix} \begin{pmatrix}
        y(0: 2 : 7) \\
        y(1: 2 : 7)
    \end{pmatrix} \\
    &= \begin{pmatrix}
        I & \Omega_4 \\
        I & -\Omega_4
    \end{pmatrix} \begin{pmatrix}
        F_4y(0: 2 : 7) \\
        F_4y(1: 2 : 7)
    \end{pmatrix}
\end{align*}

We can use the same method to calculate \begin{equation}
    F_4 y(0: 2 : 7) \mbox{ and } F_4 y(1: 2 : 7) \label{eq:F4y}
\end{equation}

Now we can have the recursive setting \begin{itemize}
    \item For $F_8y$, multiplying with $I$, $\Omega$ takes $O(m)$ time since $\Omega$ is a diagonal matrix.
    \item Then for \eqref{eq:F4y}, we can also use $O(m)$ time, which is the same as the case of $F_8y$.
    \item The number of steps is $O(\log_2 m)$.
\end{itemize}
The total time complexity is \red{$O(m \log_2 m)$}

\subsection{Matrix-product Form of FFT}

In pratice, FFT is done by \red{matrix products rather
than a recursive implementation}.

\begin{definition}[Matrix Factorization of FFT]
    Assume $m = 2^p$ for some positive integer $p$. Then we can write \[
        F = A_t \cdots A_1 P
    \]
    where \begin{gather*}
    t = \log 2m \\
    t -1 = \log m \\
    \vdots
    \end{gather*}
    and $P$ is a permutation matrix \[
        A_t = I_r \otimes B_L
    \]
    where $L = 2^{t-1},\ r = 2m / L$, and $I_r$ is the $r \times r$ identity matrix.
    
    \begin{notation}[Kronecker product]
        \[
            A \otimes B = \begin{pmatrix}
                a_{11}B & a_{12}B & \cdots & a_{1n}B \\
                \vdots & \vdots & \ddots & \vdots \\
                a_{m1}B & a_{m2}B & \cdots & a_{mn}B
            \end{pmatrix}
        \]
    \end{notation}
    
    $B_L$ is the $L \times L$ matrix defined as \[
        B_L = \begin{pmatrix}
            I_{L/2} & \Omega_{L/2} \\
            I_{L/2} & -\Omega_{L/2}
        \end{pmatrix}
    \]\[
        \Omega_{L/2} = \text{diag}(1, \delta, \delta^2, \dots, \delta^{L/2 - 1})
    \]
\end{definition}

\newpage

When $t = 3$ \[
    L = 2,\ m = 8,\ r = 1,\ I_r = 1
\]
\[
    B_8 = \begin{pmatrix}
        1& & & &1 & & & \\
        & 1 & & & & \delta & & \\  
        & & 1& & & & \delta^2 & \\
        & & & 1& & & & \delta^3 \\
        1& & & &-1 & & & \\
        & 1 & & & & -\delta & & \\  
        & & 1& & & & -\delta^2 & \\
        & & & 1& & & & -\delta^3 \\
    \end{pmatrix}
\]
When $t = 2$ \[
    L = 4,\ m = 4,\ r = 2,\ I_r = \begin{pmatrix}
        1 & 0 \\
        0 & 1
    \end{pmatrix}, B_4 = \begin{pmatrix}
        1 & 1 & 1 & 1 \\
        1 & \delta^2 & -1 & -\delta^2 \\
        1 & -1 & 1 & -1 \\
        1 & -\delta^2 & -1 & \delta^2
    \end{pmatrix}
\]
When $t = 1$ \[
    L = 8,\ m = 2,\ r = 4,\ B_2 = \begin{pmatrix}
        1 & 1 \\
        1 & -1
    \end{pmatrix}
\]
The product of $A_t \cdots A_1$ is \[
    A_3 A_2 A_1 = B_8 \begin{pmatrix}
        B_4 & 0 \\
        0 & B_4
    \end{pmatrix} \begin{pmatrix}
        B_2 & 0 & 0 & 0 \\
        0 & B_2 & 0 & 0 \\
        0 & 0 & B_2 & 0 \\
        0 & 0 & 0 & B_2
    \end{pmatrix}
\]
\[
    B_4 \begin{pmatrix}
        B_2 & 0 \\
        0 & B_2
    \end{pmatrix} = \begin{pmatrix}
        1& & 1& \\
        & 1& & \delta^2\\
        1& & -1& \\
        & 1& & -\delta^2\\
    \end{pmatrix} \begin{pmatrix}
        1&1 & & \\
        1&-1 & & \\
        & & 1&1 \\
        & & 1&-1\\
    \end{pmatrix} = \begin{pmatrix}
        1& 1& 1& 1\\  
        1& -1& \delta^2& -\delta^2\\
        1& 1& -1& -1\\
        1& -1& -\delta^2& \delta^2\\
    \end{pmatrix}
\]
Finally, we have \[
    A_3 A_2 A_1 = \begin{pmatrix}
        1 & 1 & 1 & 1 & 1 & 1 & 1 & 1 \\
        1 & -1 & \delta^2 & -\delta^2 & \delta^1 & -\delta^1 & \delta^3 & -\delta^3 \\
        1 & 1 & -1 & -1 & \delta^2 & \delta^2 & -\delta^2 & -\delta^2 \\
        1 & -1 & -\delta^2 & \delta^2 & \delta^3 & -\delta^3 & \delta^1 & -\delta^1 \\
        1 & 1 & 1 & 1 & -1 & -1 & -1 & -1 \\
        1 & -1 & \delta^2 & -\delta^2 & -\delta^1 & \delta^1 & -\delta^3 & +\delta^3 \\
        1 & 1 & -1 & -1 & -\delta^2 & -\delta^2 & \delta^2 & \delta^2 \\
        1 & -1 & -\delta^2 & \delta^2 & -\delta^3 & \delta^3 & -\delta^1 & \delta^1 \\
    \end{pmatrix}
\]
We can see that this is just the column rearrangement of $F_8$ in \eqref{eq:F4}.

\newpage

Now disscuss the permutation matrix $P$. We can see that \[
    A_3 A_2 A_1 P,
\]
where \[
    P = 
    \bordermatrix{
    & 0 & 1 & 2 & 3 & 4 & 5 & 6 & 7\cr
    0&1&&&&&&& \cr
    1&&&&&1&&& \cr
    2&&&1&&&&& \cr
    3&&&&&&&1& \cr
    4&&1&&&&&& \cr
    5&&&&&&1&& \cr
    6&&&&1&&&& \cr
    7&&&&&&&&1
    }
\]
goes back to the original $F$ \[
    F = \begin{pmatrix}
        1 & 1 & 1 & 1 & 1 & 1 & 1 & 1 \\
        1 & \delta^1 & \delta^2 & \delta^3 & -1 & -\delta^1 & -\delta^2 & -\delta^3 \\
        1 & \delta^2 & -1 & -\delta^2 & 1 & \delta^2 & -1 & -\delta^2 \\
        1 & \delta^3 & -\delta^2 & \delta & -1 & -\delta^3 & \delta^2 & -\delta^1 \\
        1 & -1 & 1 & -1 & 1 & -1 & 1 & -1 \\
        1 & -\delta^1 & \delta^2 & -\delta^3 & -1 & \delta^1 & -\delta^2 & \delta^3 \\
        1 & -\delta^2 & -1 & \delta^2 & 1 & -\delta^2 & -1 & \delta^2 \\
        1 & -\delta^3 & -\delta^2 & -\delta & -1 & \delta^3 & \delta^2 & \delta^1 \\  
    \end{pmatrix}
\]
We can see that $P$ is a permutation matrix by reversing the order of ``binary digits'' of indices.
\begin{center}
\begin{tabular}{rrrr}
    0 & 000 & 000 & 0 \\
    1 & 001 & 100 & 4 \\
    2 & 010 & 010 & 2 \\
    3 & 011 & 110 & 6 \\
    4 & 100 & 001 & 1 \\
    5 & 101 & 101 & 5 \\
    6 & 110 & 011 & 3 \\
    7 & 111 & 111 & 7
\end{tabular}
\end{center}

Futher, $A_t \cdots A_1 P$ permutes the \red{columns} of $A_t \cdots A_1$. Here \begin{align*}
    F_8 y & = \begin{pmatrix}
        I & \Omega_4 \\
        I & -\Omega_4
    \end{pmatrix} \begin{pmatrix}
        F_4 & 0 \\
        0 & F_4
    \end{pmatrix} \begin{pmatrix}
        y(0: 2 : 7) \\
        y(1: 2 : 7)
    \end{pmatrix} \\
    & = \begin{pmatrix}
        I & \Omega_4 \\
        I & -\Omega_4
    \end{pmatrix} \begin{pmatrix}
        F_4 & 0 \\
        0 & F_4
    \end{pmatrix} \begin{pmatrix}
        1   &   0   &   0   &   0   &   0   &   0   &   0   &   0   \\
        0   &   0   &   1   &   0   &   0   &   0   &   0   &   0   \\
        0   &   0   &   0   &   0   &   1   &   0   &   0   &   0   \\
        0   &   0   &   0   &   0   &   0   &   0   &   1   &   0   \\
        0   &   1   &   0   &   0   &   0   &   0   &   0   &   0   \\
        0   &   0   &   0   &   1   &   0   &   0   &   0   &   0   \\
        0   &   0   &   0   &   0   &   0   &   1   &   0   &   0   \\
        0   &   0   &   0   &   0   &   0   &   0   &   0   &   1   \\
    \end{pmatrix}
    y
\end{align*}

\newpage

We can see that $y$ is rearranged according to the following order: \begin{center}
\begin{tabular}{rrrrr}
    0 & 000 &               & 0 & 000  \\
    1 & 001 &               & 2 & 010  \\
    2 & 010 &               & 4 & 100  \\
    3 & 011 &$\Rightarrow$  & 6 & 110  \\
    4 & 100 &               & 1 & 001  \\
    5 & 101 &               & 3 & 011  \\
    6 & 110 &               & 5 & 101  \\
    7 & 111 &               & 7 & 111 
\end{tabular}
\end{center}

This is like that the binary representation is \red{rotated to the left by one digit}.

For example,
\begin{equation*}
    5 \Rightarrow 2(5-4) + 1 = 3
\end{equation*}
\begin{tabular}{ll}
    $5-4$:  &First digit ``1'' removed\\
    $2(5-4)$: &  Second, third digits shifted left\\
    $2(5-4)+1$: & Add the original first digit ``1'' to the now last digit
\end{tabular}

Now we can see that besides the last digit, the \red{first two} digits of both the upper and the lower parts are the same:
\begin{center}
\begin{tabular}{rrrr}
    0 & 00 \\
    1 & 01 \\
    2 & 10 \\
    3 & 11 \\
\end{tabular}
\end{center}
It is like we have a new small FFT problem \[
    F_4 \tilde{y}
\]
By the same setting we should have \[
    F_8 y  = A_3 \begin{pmatrix}
    I_2     &\Omega_2   &   0   &   0\\
    I_2     &-\Omega_2  &   0   &   0\\
    0       &   0       &   I_2 &   \Omega_2\\
    0       &   0       &   I_2 &   -\Omega_2
\end{pmatrix}
\begin{pmatrix}
    F_2     &   0   &   0   &   0   \\
    0       &   F_2 &   0   &   0   \\
    0       &   0   &   F_2 &   0   \\
    0       &   0   &   0   &   F_2 \\
    \end{pmatrix}
    \underbrace{  \begin{pmatrix}
        1   &   0   &   0   &   0   &   0   &   0   &   0   &   0   \\
        0   &   0   &   1   &   0   &   0   &   0   &   0   &   0   \\
        0   &   1   &   0   &   0   &   0   &   0   &   0   &   0   \\
        0   &   0   &   0   &   1   &   0   &   0   &   0   &   0   \\
        0   &   0   &   0   &   0   &   1   &   0   &   0   &   0   \\
        0   &   0   &   0   &   0   &   0   &   0   &   1   &   0   \\
        0   &   0   &   0   &   0   &   0   &   1   &   0   &   0   \\
        0   &   0   &   0   &   0   &   0   &   0   &   0   &   1   \\
        \end{pmatrix}
}_{\text{block diagonal}}                      
    \begin{pmatrix}
        1   &   0   &   0   &   0   &   0   &   0   &   0   &   0   \\
        0   &   0   &   1   &   0   &   0   &   0   &   0   &   0   \\
        0   &   0   &   0   &   0   &   1   &   0   &   0   &   0   \\
        0   &   0   &   0   &   0   &   0   &   0   &   1   &   0   \\
        0   &   1   &   0   &   0   &   0   &   0   &   0   &   0   \\
        0   &   0   &   0   &   1   &   0   &   0   &   0   &   0   \\
        0   &   0   &   0   &   0   &   0   &   1   &   0   &   0   \\
        0   &   0   &   0   &   0   &   0   &   0   &   0   &   1   \\
    \end{pmatrix}y
\]
We have \[
    A_2 = \begin{pmatrix}
        I_2     &\Omega_2   &   0   &   0\\
        I_2     &-\Omega_2  &   0   &   0\\
        0       &   0       &   I_2 &   \Omega_2\\
        0       &   0       &   I_2 &   -\Omega_2
    \end{pmatrix}
\]
For the next step, we have the permutation matrix is $I$. Thus \[
    P = I_8 \begin{pmatrix}
        1   &   0   &   0   &   0   &   0   &   0   &   0   &   0   \\
        0   &   0   &   1   &   0   &   0   &   0   &   0   &   0   \\
        0   &   1   &   0   &   0   &   0   &   0   &   0   &   0   \\
        0   &   0   &   0   &   1   &   0   &   0   &   0   &   0   \\
        0   &   0   &   0   &   0   &   1   &   0   &   0   &   0   \\
        0   &   0   &   0   &   0   &   0   &   0   &   1   &   0   \\
        0   &   0   &   0   &   0   &   0   &   1   &   0   &   0   \\
        0   &   0   &   0   &   0   &   0   &   0   &   0   &   1   \\
    \end{pmatrix}                     
    \begin{pmatrix}
        1   &   0   &   0   &   0   &   0   &   0   &   0   &   0   \\
        0   &   0   &   1   &   0   &   0   &   0   &   0   &   0   \\
        0   &   0   &   0   &   0   &   1   &   0   &   0   &   0   \\
        0   &   0   &   0   &   0   &   0   &   0   &   1   &   0   \\
        0   &   1   &   0   &   0   &   0   &   0   &   0   &   0   \\
        0   &   0   &   0   &   1   &   0   &   0   &   0   &   0   \\
        0   &   0   &   0   &   0   &   0   &   1   &   0   &   0   \\
        0   &   0   &   0   &   0   &   0   &   0   &   0   &   1   \\
    \end{pmatrix} = \begin{pmatrix}
        1   &   0   &   0   &   0   &   0   &   0   &   0   &   0\\
        0   &   0   &   0   &   0   &   1   &   0   &   0   &   0\\
        0   &   0   &   1   &   0   &   0   &   0   &   0   &   0\\
        0   &   0   &   0   &   0   &   0   &   0   &   1   &   0\\
        0   &   1   &   0   &   0   &   0   &   0   &   0   &   0\\
        0   &   0   &   0   &   0   &   0   &   1   &   0   &   0\\
        0   &   0   &   0   &   1   &   0   &   0   &   0   &   0\\
        0   &   0   &   0   &   0   &   0   &   0   &   0   &   1\\
    \end{pmatrix}
\]
and \[
    A_1 = \begin{pmatrix}
        I_1     &\Omega_1   &   0   &   0   &   0   &   0   &   0   &   0\\
        I_1     &-\Omega_1  &   0   &   0   &   0   &   0   &   0   &   0\\
        0       &   0       &I_1     &\Omega_1 &   0   &   0   &   0   &   0\\
        0       &   0       &I_1     &-\Omega_1&   0   &   0   &   0   &   0\\
        0       &   0       &0       &   0       &I_1     &\Omega_1 &   0   &   0\\
        0       &   0       &0       &   0       &I_1     &-\Omega_1&   0   &   0\\
        0       &   0       &0       &   0       &0       &   0       &I_1     &\Omega_1 \\
        0       &   0       &0       &   0       &0       &   0       &I_1     &-\Omega_1\\
    \end{pmatrix}
\]
Therefore, $y$ is arranged by the following way
\begin{center}
\begin{tabular}{rrrrrrrr}
    0 & 000 &               & 0 & 000 &             & 0& 000     \\
    1 & 001 &               & 2 & 010 &             & 4& 100     \\
    2 & 010 &               & 4 & 100 &             & 2& 010     \\
    3 & 011 & $\Rightarrow$ & 6 & 110 &$\Rightarrow$& 6& 110     \\
    4 & 100 &               & 1 & 001 &             & 1& 001     \\
    5 & 101 &               & 3 & 011 &             & 5& 101     \\
    6 & 110 &               & 5 & 101 &             & 3& 011     \\
    7 & 111 &               & 7 & 111 &             & 7& 111     
\end{tabular}
\end{center}
Thus we have explained the reason of using the reverse of the binary digits.